\documentclass[10.5pt]{config}
% 本实验报告采用 署名-相同方式共享 4.0 国际许可协议 (CC BY-SA 4.0)
\usepackage{amssymb}   % 提供 \mathbb 等数学符号

\major{电子科学与技术}
\name{郑智瑄}
\partner{} % 签名栏
\stuid{3230103913}
\date{\today}
\lab{玉泉高分子楼-203}
\course{人工智能实验}
\instructor{魏准}
\grades{}
\expname{Linear / Nonlinear SVM} % 实验名字
\exptype{设计型实验} % 实验类型

% minted 代码样式设置
\usemintedstyle{tango}
\setminted{
    fontsize=\small,
    linenos,
    numbersep=8pt,
    breaklines,
    frame=single,
    framesep=2mm,
    tabsize=4,
    bgcolor=codebg
}
% 定义代码背景色
\definecolor{codebg}{RGB}{248, 248, 248}

\begin{document}

\makeheader % 首页头部

% ==================== SVM 线性分类 ====================
\section{SVM 线性分类}

\subsection{实验目的}
基于 Scikit-Learn 的支持向量机模块，使用鸢尾花数据集中的花瓣长度与宽度两个特征，完成线性 SVM 二分类（维吉尼亚鸢尾花 vs 其余类别）；重点测试花瓣长度 $\in[4,6]$、花瓣宽度 $\in[0.8,2.5]$ 监测区域内各点的分类情况，可视化决策超平面、间隔边界及支持向量，并对比不同 $C$ 值的影响。

\subsection{实验原理}
支持向量机（SVM）的目标是在特征空间中找到使两类样本间隔最大的超平面。对于线性可分情形，最优决策超平面满足：
$$\mathbf{w}^{\top}\mathbf{x} + b = 0, \quad \text{间隔宽度} = \frac{2}{\|\mathbf{w}\|}$$
为获得最大间隔，等价于求解以下凸优化问题：
$$\min_{\mathbf{w},b}\ \frac{1}{2}\|\mathbf{w}\|^2 \quad \text{s.t.}\ y_i(\mathbf{w}^{\top}\mathbf{x}_i+b)\geq 1,\ \forall i$$
软间隔 SVM 引入松弛变量 $\xi_i \geq 0$ 和正则化超参数 $C$，允许少量样本落在间隔内部：
$$\min_{\mathbf{w},b,\boldsymbol{\xi}}\ \frac{1}{2}\|\mathbf{w}\|^2 + C\sum_i\xi_i \quad \text{s.t.}\ y_i(\mathbf{w}^{\top}\mathbf{x}_i+b)\geq 1-\xi_i,\ \xi_i\geq 0$$
决策超平面由位于间隔边界上（或内部）的少数样本（支持向量）完全决定。

\subsubsection{Pipeline 工作流}
实验代码使用 Scikit-Learn 的 \texttt{Pipeline} 将数据预处理与模型训练串联为一个整体，确保训练与预测阶段使用完全相同的变换参数：
$$\text{原始特征} \xrightarrow{\texttt{StandardScaler}} \text{标准化特征（零均值、单位方差）} \xrightarrow{\texttt{SVC(kernel="linear", C)}} \text{分类结果}$$
其中 \texttt{StandardScaler} 对花瓣长度和宽度进行标准化，避免量纲差异影响 SVM 的间隔计算；\texttt{SVC} 以线性核在标准化特征空间中求解最优超平面。

\subsection{核心代码实现}
\begin{listing}[H]
\begin{minted}{python}
from sklearn import datasets
from sklearn.svm import SVC
from sklearn.preprocessing import StandardScaler
from sklearn.pipeline import Pipeline
import numpy as np

# ① 加载数据集
iris = datasets.load_iris()
X = iris["data"][:, (2, 3)]            # 花瓣长度、花瓣宽度
y = (iris["target"] == 2).astype(int)  # 1: Virginica, 0: others

# ② 缩放特征 + ③ 训练线性 SVM
svm_clf = Pipeline([
    ("scaler", StandardScaler()),
    ("svc", SVC(kernel="linear", C=1.0, random_state=42)),
])
svm_clf.fit(X, y)

# ④ 监测区域：花瓣长度 4~6 cm，花瓣宽度 0.8~2.5 cm
pl_vals = np.arange(4.0, 6.1, 0.2)
pw_vals = np.arange(0.8, 2.6, 0.2)
pl_grid, pw_grid = np.meshgrid(pl_vals, pw_vals)
monitor_pts  = np.c_[pl_grid.ravel(), pw_grid.ravel()]
monitor_pred = svm_clf.predict(monitor_pts)

n_virginica = int(monitor_pred.sum())
n_other     = len(monitor_pred) - n_virginica
print(f"监测点总数: {len(monitor_pred)}")
print(f"判定为 Virginica    : {n_virginica}")
print(f"判定为 Non-Virginica: {n_other}")
\end{minted}
\caption{SVM 线性分类核心代码（节选）}
\end{listing}

\subsection{运行结果}
程序输出如下：
\begin{lstlisting}[language={}, numbers=none, frame=single, backgroundcolor=\color{codebg}, basicstyle=\ttfamily\footnotesize, breaklines=true]
监测点总数         : 99
判定为 Virginica   : 48 个
判定为 Non-Virginica: 51 个

  花瓣长度  花瓣宽度       预测类别    决策值
  --------  --------  ----------  --------
       ...       ...           ...       ...
       5.6       1.4  Non-Virginica   -0.0126  <- 临界
       5.8       1.4      Virginica    0.2254  <- 越界
       4.0       2.0      Virginica    0.0126  <- 越界
       ...       ...           ...       ...

Training accuracy : 0.9533
Number of support vectors: [13 13] (class 0 / class 1)
Decision function (scaled): 2.0935*x1 + 2.4425*x2 + -2.8496 = 0
Margin width (scaled)     : 0.6217
\end{lstlisting}

\begin{figure}[H]
    \centering
    \includegraphics[width=0.82\textwidth]{Figure_7.png}
    \caption{SVM 线性分类结果；紫色点划线矩形为监测区域 $[4,6]\times[0.8,2.5]$，×号为监测点预测结果（蓝×为 Non-Virginica，红×为 Virginica）}
    \label{fig:svm_linear}
\end{figure}

\subsection{结果分析}
如图~\ref{fig:svm_linear} 所示，线性 SVM 在标准化特征空间中找到了一条明确的分类超平面（黑色实线），并以虚线标出两侧等间隔边界。圆圈标注的 26 个支持向量（每类 13 个）完全决定了该超平面的位置与方向。训练集准确率为 0.9533，模型已将绝大多数维吉尼亚鸢尾花与其他类别正确区分。

对于监测区域（花瓣长度 $\in[4,6]$，花瓣宽度 $\in[0.8,2.5]$，共 99 个采样点）：
\begin{itemize}
    \item \textbf{花瓣宽度 $<1.4$ cm} 的所有点：决策值均为负，全部判定为 Non-Virginica；
    \item \textbf{花瓣宽度 $\geq 2.0$ cm} 的所有点：决策值均为正，全部判定为 Virginica；
    \item \textbf{花瓣宽度 $\in[1.4,1.8]$ cm} 区域（约占 1/3）：为决策边界过渡带，分类结果同时取决于花瓣长度——花瓣越长、宽度越大越倾向于 Virginica。最终 99 个监测点中有 48 个判为 Virginica（48.5\%），51 个判为 Non-Virginica（51.5\%）。
\end{itemize}
这与决策超平面 $2.09\,x_1 + 2.44\,x_2 - 2.85 = 0$ 的数学形式一致：\textbf{花瓣宽度对分类的权重（2.44）略大于花瓣长度（2.09）}，说明在监测区域中花瓣宽度是更主要的判别依据。

\subsection{不同 \texorpdfstring{$C$}{C} 值对分类结果的影响}

\subsubsection{实验设置}
保持线性核不变，将正则化超参数 $C$ 依次取 0.5、1、15、30、50、100，观察间隔宽度、支持向量数量及训练准确率的变化规律。核心代码如下：

\begin{listing}[H]
\begin{minted}{python}
C_values = [0.5, 1, 15, 30, 50, 100]
for C in C_values:
    svm_clf = Pipeline([
        ("scaler", StandardScaler()),
        ("svc", SVC(kernel="linear", C=C, random_state=42)),
    ])
    svm_clf.fit(X, y)
    svc    = svm_clf.named_steps["svc"]
    acc    = svm_clf.score(X, y)
    margin = 2 / np.linalg.norm(svc.coef_[0])
    n_sv   = svc.n_support_
    print(f"C={C:6.1f}  acc={acc:.4f}  margin={margin:.4f}  n_sv={n_sv}")
\end{minted}
\caption{SVM 线性分类——不同 $C$ 值比较核心代码}
\end{listing}

\subsubsection{运行结果}
程序输出如下：
\begin{lstlisting}[language={}, numbers=none, frame=single, backgroundcolor=\color{codebg}, basicstyle=\ttfamily\footnotesize, breaklines=true]
     C    Accuracy      Margin      n_sv (0/1)
--------------------------------------------------
   0.5      0.9533      0.6954     16 / 16
   1.0      0.9533      0.6217     13 / 13
  15.0      0.9667      0.3409      7 / 7
  30.0      0.9667      0.3288      7 / 7
  50.0      0.9600      0.2494      6 / 7
 100.0      0.9600      0.2154      6 / 6
\end{lstlisting}

六种 $C$ 值的决策边界可视化见图~\ref{fig:svm_c_compare}。

\begin{figure}[H]
    \centering
    \includegraphics[width=0.9\textwidth]{Figure_10.png}
    \caption{SVM 线性分类——不同 $C$ 值下的决策边界（实线为决策面，虚线为间隔边界，圆圈为支持向量）}
    \label{fig:svm_c_compare}
\end{figure}

\subsubsection{结果分析与结论}

\paragraph{$C$ 值与间隔宽度的关系}
$C$ 是软间隔 SVM 中惩罚误分类样本的正则化系数。
\begin{itemize}
    \item \textbf{$C$ 较小（0.5、1）}：惩罚力度弱，模型优先追求更宽的间隔，允许更多样本落入间隔内甚至被误分类。间隔宽度分别为 0.6954 和 0.6217，支持向量数多（各 16 个、13 个），决策超平面位置由较多样本共同决定，方向更"平均"。
    \item \textbf{$C$ 适中（15、30）}：惩罚力度增强，模型在减小训练误差与保持间隔之间取得平衡，准确率提升至 0.9667，间隔压缩至约 0.34，支持向量减少至 7 个/类。
    \item \textbf{$C$ 较大（50、100）}：模型进一步减少容忍的误分类样本，但准确率反而略降至 0.9600，间隔收窄至 0.25 和 0.22，支持向量最少（约 6 个/类）。这说明 $C$ 过大时模型开始对噪声点过度拟合，泛化能力略有下降。
\end{itemize}

\paragraph{量化汇总}
\begin{table}[H]
    \centering
    \caption{不同 $C$ 值下的关键指标对比}
    \begin{tabular}{ccccc}
        \hline
        $C$ & 训练准确率 & 间隔宽度（缩放空间） & 支持向量数（0类/1类） \\
        \hline
        0.5  & 0.9533 & 0.6954 & 16 / 16 \\
        1    & 0.9533 & 0.6217 & 13 / 13 \\
        15   & 0.9667 & 0.3409 &  7 /  7 \\
        30   & 0.9667 & 0.3288 &  7 /  7 \\
        50   & 0.9600 & 0.2494 &  6 /  7 \\
        100  & 0.9600 & 0.2154 &  6 /  6 \\
        \hline
    \end{tabular}
    \label{tab:c_compare}
\end{table}

\paragraph{结论}
$C$ 值通过调节误分类惩罚力度，在\textbf{间隔宽度}（泛化能力）与\textbf{训练误差}（拟合能力）之间形成权衡：
\begin{enumerate}
    \item $C$ 增大 $\Rightarrow$ 间隔变窄、支持向量减少、训练误差降低，但过大时有过拟合风险；
    \item $C$ 减小 $\Rightarrow$ 间隔变宽、支持向量增多、模型更"保守"，泛化能力通常更强；
    \item 本数据集上 $C\in[15,30]$ 取得了准确率与间隔的最佳平衡（准确率 0.9667，间隔约 0.34）；$C$ 继续增大至 50、100 后训练准确率不升反降，说明已出现轻微过拟合。
    \item 实践中应通过\textbf{交叉验证}（如 \texttt{GridSearchCV}）在保留验证集上选取最优 $C$，而非单纯依赖训练集准确率。
\end{enumerate}

% ==================== SVM 非线性分类 I ====================
\clearpage
\section{SVM 非线性分类 I（多项式特征 + LinearSVC）}

\subsection{实验目的}
使用 Scikit-Learn 的线性分类器与多项式特征，对 \texttt{make\_moons} 生成的月牙形非线性数据集完成二分类，具体目标如下：
\begin{enumerate}
    \item 产生非线性数据集（\texttt{make\_moons}）并划分训练集与测试集；
    \item 利用 \texttt{Pipeline} 依次组合 \texttt{PolynomialFeatures}、\texttt{StandardScaler} 与 \texttt{LinearSVC}，构建非线性分类器；
    \item 可视化训练数据与测试数据；
    \item 讨论分析不同多项式阶数（2 阶、3 阶、8 阶）对分类结果的影响；
    \item 讨论分析不同正则化参数（$C=1,10,100$）对分类结果的影响。
\end{enumerate}

\subsection{实验原理}

\subsubsection{非线性 SVM 的基本思路}
对于线性不可分数据集，一种直观的解决方案是将原始特征空间 $\mathbf{x}\in\mathbb{R}^p$ 通过显式的\textbf{多项式特征拓展}映射到高维空间 $\tilde{\mathbf{x}}\in\mathbb{R}^{p'}$，使得数据在新空间中线性可分，再用线性 SVM 进行分类：
$$\tilde{\mathbf{x}} = (X_1,\ X_1^2,\ X_2,\ X_2^2,\ \ldots,\ X_p,\ X_p^2,\ X_1 X_2,\ \ldots)$$
超平面在原始特征空间中表现为非线性决策边界：
$$\beta_0 + \beta_1 X_1 + \beta_2 X_1^2 + \cdots + \beta_{2p-1} X_p + \beta_{2p} X_p^2 = 0$$

\subsubsection{软间隔线性 SVM（LinearSVC）}
软间隔 SVM 引入松弛变量 $\xi_i \geq 0$ 与惩罚系数 $C$，在拓展的特征空间上求解：
$$\min_{\mathbf{w},b,\boldsymbol{\xi}}\ \frac{1}{2}\|\mathbf{w}\|^2 + C\sum_{i=1}^{n}\xi_i$$
$$\text{s.t.}\quad y_i(\mathbf{w}^{\top}\tilde{\mathbf{x}}_i + b) \geq 1-\xi_i,\quad \xi_i\geq 0$$
其中 $C$ 越大，对误分类的惩罚越重，决策边界越贴近训练样本（易过拟合）；$C$ 越小，间隔越宽，泛化能力通常更强。\texttt{LinearSVC} 采用 hinge 损失函数，等价于上述软间隔优化问题的对偶形式。

\subsubsection{Pipeline 工作流}
\texttt{Pipeline} 将数据预处理与模型训练串联为一个整体：
$$\text{原始特征} \xrightarrow{\text{PolynomialFeatures(degree)}} \text{高维特征} \xrightarrow{\text{StandardScaler}} \text{标准化} \xrightarrow{\text{LinearSVC(C)}} \text{分类结果}$$
标准化（零均值、单位方差）在多项式拓展后执行，避免高次项数值量级过大影响优化稳定性。

\subsection{核心代码实现}
\begin{listing}[H]
\begin{minted}{python}
import numpy as np
import matplotlib.pyplot as plt
from sklearn.datasets import make_moons
from sklearn.pipeline import Pipeline
from sklearn.preprocessing import StandardScaler, PolynomialFeatures
from sklearn.svm import LinearSVC
from sklearn.model_selection import train_test_split
from sklearn.metrics import classification_report, confusion_matrix

# 1. 生成月牙形数据集并划分训练/测试集
X, y = make_moons(n_samples=200, noise=0.15, random_state=42)
X_train, X_test, y_train, y_test = train_test_split(
    X, y, test_size=0.25, random_state=42   # 150 train / 50 test
)

# 2. 构建 Pipeline：PolynomialFeatures + StandardScaler + LinearSVC
polynomial_svm_clf = Pipeline([
    ("poly_features", PolynomialFeatures(degree=3)),   # 产生非线性特征
    ("scaler",        StandardScaler()),                # 标准化
    ("svm_clf",       LinearSVC(C=10, loss="hinge",   # 线性 SVM (hinge 损失)
                                max_iter=5000, random_state=42)),
])

# 3. 训练与预测
polynomial_svm_clf.fit(X_train, y_train)
y_pred = polynomial_svm_clf.predict(X_test)

# 4. 评估
print(f"Train accuracy: {polynomial_svm_clf.score(X_train, y_train):.4f}")
print(f"Test  accuracy: {polynomial_svm_clf.score(X_test,  y_test ):.4f}")
print(classification_report(y_test, y_pred))
print("Confusion matrix:\n", confusion_matrix(y_test, y_pred))

# 5. 不同阶数对比（C=1 固定）
for deg in [2, 3, 8]:
    clf = Pipeline([
        ("poly_features", PolynomialFeatures(degree=deg)),
        ("scaler",        StandardScaler()),
        ("svm_clf",       LinearSVC(C=1, loss="hinge",
                                    max_iter=5000, random_state=42)),
    ])
    clf.fit(X_train, y_train)
    print(f"Degree={deg}  Train={clf.score(X_train,y_train):.4f}"
          f"  Test={clf.score(X_test,y_test):.4f}")

# 6. 不同 C 值对比（degree=3 固定）
for C in [1, 10, 100]:
    clf = Pipeline([
        ("poly_features", PolynomialFeatures(degree=3)),
        ("scaler",        StandardScaler()),
        ("svm_clf",       LinearSVC(C=C, loss="hinge",
                                    max_iter=5000, random_state=42)),
    ])
    clf.fit(X_train, y_train)
    print(f"C={C:<4}    Train={clf.score(X_train,y_train):.4f}"
          f"  Test={clf.score(X_test,y_test):.4f}")
\end{minted}
\caption{SVM 非线性分类 I 核心代码}
\end{listing}

\subsection{数据可视化}

\begin{figure}[H]
    \centering
    \includegraphics[width=0.55\textwidth]{svm_poly_data.png}
    \caption{make\_moons 数据集可视化：圆点为训练集（80 个样本），星号为测试集（20 个样本），蓝色为 Class~0，红色为 Class~1}
    \label{fig:svm_poly_data}
\end{figure}

如图~\ref{fig:svm_poly_data} 所示，月牙形数据集由两条相互缠绕的弧形点云组成，两类样本在原始二维特征空间中明显线性不可分，任何线性超平面均无法将其正确分开，需借助多项式特征拓展至高维空间。数据集共 200 个样本，按 3:1 划分训练集（150 个）和测试集（50 个），步长为 0.02，能够更真实地反映模型泛化能力。

\subsection{运行结果}

\subsubsection{控制台输出}
\begin{lstlisting}[language={}, numbers=none, frame=single, backgroundcolor=\color{codebg}, basicstyle=\ttfamily\footnotesize, breaklines=true]
训练集大小: 150  测试集大小: 50

【实验A】不同多项式阶数 (C=1)
  Degree   Train Acc    Test Acc
  ---------------------------------
       2      0.8533      0.8600
       3      0.9667      0.9800
       8      0.9800      0.9800

【实验B】不同 C 值 (degree=3)
       C   Train Acc    Test Acc
  ---------------------------------
       1      0.9667      0.9800
      10      1.0000      0.9800
     100      1.0000      0.9800

【最优模型详细报告】degree=3, C=10
              precision    recall  f1-score   support
           0       0.96      1.00      0.98        26
           1       1.00      0.96      0.98        24
    accuracy                           0.98        50
   macro avg       0.98      0.98      0.98        50
weighted avg       0.98      0.98      0.98        50

混淆矩阵:
[[26  0]
 [ 1 23]]
\end{lstlisting}

\subsubsection{不同多项式阶数的决策边界}

\begin{figure}[H]
    \centering
    \includegraphics[width=0.98\textwidth]{svm_poly_degree.png}
    \caption{不同多项式阶数（$C=1$ 固定）的决策边界；圆点为训练集，星号为测试集；背景色为预测区域，黑线为决策边界}
    \label{fig:svm_poly_degree}
\end{figure}

\subsubsection{不同 $C$ 值的决策边界}

\begin{figure}[H]
    \centering
    \includegraphics[width=0.98\textwidth]{svm_poly_C.png}
    \caption{不同正则化系数 $C$（Degree=3 固定）的决策边界；随 $C$ 增大，决策边界更贴近训练点，过拟合风险升高}
    \label{fig:svm_poly_C}
\end{figure}

\subsection{结果分析}

\subsubsection{不同多项式阶数的影响}

\begin{table}[H]
    \centering
    \caption{不同多项式阶数对比（$C=1$ 固定，训练集 150 / 测试集 50）}
    \begin{tabular}{cccc}
        \hline
        阶数 (Degree) & 特征维数 & 训练准确率 & 测试准确率 \\
        \hline
        2 & 6  & 0.8533 & 0.8600 \\
        3 & 10 & 0.9667 & \textbf{0.9800} \\
        8 & 45 & 0.9800 & \textbf{0.9800} \\
        \hline
    \end{tabular}
    \label{tab:degree_compare}
\end{table}

从表~\ref{tab:degree_compare} 和图~\ref{fig:svm_poly_degree} 可以得出以下结论：

\begin{itemize}
    \item \textbf{2 阶}：特征维数仅 6，模型复杂度最低。训练准确率 0.853，测试准确率 0.860，两者均为三种阶数中最低，表明 2 阶多项式边界过于简单，\textbf{欠拟合}明显——决策边界无法充分拟合月牙形曲线，仍有较多样本被误分类。
    \item \textbf{3 阶}：特征维数 10，训练与测试准确率达到 0.967 / 0.980，两者接近且均衡，是三者中综合最优的选择。决策边界足够弯曲以包络月牙形轮廓，同时未对噪声点过度敏感。
    \item \textbf{8 阶}：特征维数增至 45，训练准确率 0.980，测试准确率同为 0.980。虽然数值上与 3 阶相当，但由图可见决策边界已出现局部不规则波动（对部分噪声点过度拟合），模型的\textbf{过拟合风险}更高，且计算开销显著增大，并不优于 3 阶。
\end{itemize}

\textbf{结论}：多项式阶数控制模型复杂度。2 阶欠拟合，决策边界无法表达月牙形；3 阶在拟合能力与泛化之间取得最佳平衡；8 阶虽测试准确率与 3 阶持平，但边界更复杂、计算成本更高，实际中应优先选用 3 阶。

\subsubsection{不同 $C$ 值的影响}

\begin{table}[H]
    \centering
    \caption{不同 $C$ 值对比（Degree=3 固定，训练集 150 / 测试集 50）}
    \begin{tabular}{cccc}
        \hline
        $C$ & 训练准确率 & 测试准确率 & 特点 \\
        \hline
        1   & 0.9667 & \textbf{0.9800} & 间隔较宽，容忍少量误分类 \\
        10  & 1.0000 & \textbf{0.9800} & 训练完全正确，间隔收紧 \\
        100 & 1.0000 & \textbf{0.9800} & 间隔更窄，过拟合风险升高 \\
        \hline
    \end{tabular}
    \label{tab:C_compare}
\end{table}

从表~\ref{tab:C_compare} 和图~\ref{fig:svm_poly_C} 可以得出以下结论：

\begin{itemize}
    \item \textbf{$C=1$}：惩罚力度较弱，模型保持较宽间隔，训练准确率 0.967，测试准确率 0.980。在三个 $C$ 值中，$C=1$ 以最低的训练拟合度换得了与高 $C$ 值相同的测试表现，模型最为"保守"且泛化稳定。
    \item \textbf{$C=10$}：惩罚力度增大，训练准确率提升至 1.000（训练集完全正确），测试准确率维持 0.980，与 $C=1$ 持平。间隔收紧，决策边界更贴近训练样本，但在此数据集上尚未出现泛化下降。
    \item \textbf{$C=100$}：惩罚力度最强，训练准确率同为 1.000，测试准确率仍为 0.980，但由图可见决策边界局部出现更多不规则弯折，对噪声点的敏感程度进一步升高，\textbf{过拟合风险}最大，在更复杂或更嘈杂的数据集上泛化能力可能下降。
\end{itemize}

\textbf{结论}：在当前 Degree=3 设置下，三个 $C$ 值的测试准确率均为 0.980，差异体现在训练准确率和边界复杂度上。$C=1$ 以最宽间隔取得与更高 $C$ 值相同的泛化效果，是本实验中最优的 $C$ 选取；随 $C$ 增大，模型逐渐从"宽间隔容错"转向"零训练误差强拟合"，虽暂未导致测试准确率下降，但过拟合风险持续积累。

\clearpage
\section{实验总结}
通过两组 SVM 实验，深入理解了支持向量机的核心原理与实践方法。

\textbf{SVM 线性分类实验}验证了线性 SVM 通过最大化间隔寻找最优决策超平面的思路，并通过对 $C\in\{0.5,1,15,30,50,100\}$ 的系统比较，直观展示了软间隔参数 $C$ 在容忍误分类与增大间隔之间的权衡——$C=15\sim30$ 时准确率最高（0.9667），过大的 $C$（50、100）反而因过拟合导致准确率略降，说明选取合适的 $C$ 对模型泛化至关重要。

\textbf{SVM 非线性分类 I 实验}通过多项式特征拓展（\texttt{PolynomialFeatures}）将线性不可分的月牙形数据映射到高维空间，再使用 \texttt{LinearSVC} 线性分类，验证了"显式特征映射"实现非线性分类的可行性。实验结果表明：
\begin{enumerate}
    \item \textbf{阶数}：2 阶明显欠拟合（训练/测试均约 0.86），3 阶与 8 阶测试准确率相当（均为 0.980），但 3 阶特征维数更低（10 维 vs 45 维）、边界更简洁，为最优选择；
    \item \textbf{C 值}：Degree=3 下，$C=1,10,100$ 的测试准确率均为 0.980，但 $C=1$ 以最宽间隔取得相同泛化效果，随 $C$ 增大训练准确率达到 1.000 但过拟合风险持续积累，体现了偏差—方差权衡的核心规律；
    \item \textbf{实践意义}：多项式特征数随阶数指数增长（2 阶 6 维 $\to$ 8 阶 45 维），高阶时计算成本与过拟合风险同步上升，实际应用中应结合交叉验证选择合适阶数与 $C$ 值。
\end{enumerate}

整体而言，SVM 的核心超参数（$C$、特征阶数）对模型性能影响显著，实践中应结合交叉验证进行系统调优。SVM 算法理论完备（最大间隔与结构风险最小化），在小样本场景下具有较强的竞争力。

\end{document}
